\documentclass[12pt,a4paper]{article}

% Pacotes essenciais para português e formatação
\usepackage[utf8]{inputenc}
\usepackage[T1]{fontenc}
\usepackage[portuguese]{babel}
\usepackage{geometry}
\geometry{a4paper, margin=2.5cm}
\usepackage{amsmath}
\usepackage{amssymb}
\usepackage{verbatim} % Para os blocos de pseudocódigo
\usepackage{hyperref}

\title{Trabalho 1: Busca Adversarial em Jogos\\Agente Inteligente para Breakthrough}
\author{Pedro Henrique dos Reis Arcoverde - 254719 \and Nome do Aluno 2 \and Nome do Aluno 3 \and Nome do Aluno 4}
\date{Março de 2025}

\begin{document}

\maketitle

\section{Modelagem do Problema}

\subsection{Apresentação do Problema: Breakthrough}
O Breakthrough é um jogo de tabuleiro estratégico de soma zero e informação perfeita. 
O jogo ocorre em um tabuleiro quadriculado, onde cada jogador começa com duas fileiras de peças. 
As regras de movimentação ditam que as peças movem-se apenas para frente, ou na diagonal para capturar 
peças inimigas. O jogo possui duas condições de vitória: vence quem alcançar a última fileira adversária 
ou eliminar todas as peças do oponente.

\subsection{Definição Formal do Jogo}
Para modelar o Breakthrough como um problema de busca em Inteligência Artificial, 
definimos formalmente o jogo como a tupla $(S, A, T, U)$:

\begin{itemize}
    \item \textbf{Estados ($S$):} O conjunto de todas as configurações possíveis do tabuleiro. Cada estado é representado 
    por uma matriz $8\times8$, onde cada célula pode conter: vazio (0), peça do Jogador 1 (1) ou peça do Jogador 2 (2).
    \item \textbf{Ações ($A$):} As jogadas legais disponíveis para um jogador em um dado estado.
     Uma ação é representada pela coordenada de origem da peça e o seu destino.
    \item \textbf{Modelo de Transição ($T$):} A função $T(s, a) = s'$, que recebe um estado $s$ e uma ação legal $a$, 
    e retorna um novo estado $s'$.
    \item \textbf{Função de Utilidade ($U$):} A função que avalia os estados. Retorna um valor numérico máximo 
    se o estado resulta em vitória para o agente, um valor mínimo se resulta em derrota, e avaliações heurísticas 
    intermediárias para estados não terminais.
\end{itemize}

\subsection{Teste de Terminal}
O teste de estado terminal verifica as condições de fim de jogo. A busca é interrompida naquele ramo se houver invasão 
(uma peça alcançou a fileira base do inimigo) ou aniquilação (um jogador perdeu todas as peças).

\section{Algoritmos de Busca Utilizados}

Para dotar o agente da capacidade de tomar decisões competitivas, foram implementados algoritmos clássicos 
de busca adversarial, respeitando o limite orçamentário de tempo por jogada e utilizando \textit{Iterative Deepening}.

\subsection{Minimax Puro}
O algoritmo Minimax assume que ambos os jogadores jogam de forma otimizada. O agente (Max) busca maximizar 
sua função de utilidade, enquanto o oponente (Min) busca minimizá-la. Seu custo computacional é $\mathcal{O}(b^d)$, 
onde $b$ é o fator de ramificação e $d$ é a profundidade. Sem podas, o Minimax puro se comporta fazendo uma busca 
exaustiva em cada ramificação da árvore do jogo o que acaba sendo inviável para profundidades maiores devido à explosão combinatória, como
também, não é funcional dados os limites de tempo impostos por jogada.

\begin{verbatim}
Função MINIMAX(estado, profundidade, maximizando):
    Se profundidade == 0 ou TESTE-TERMINAL(estado):
        Retorna AVALIAÇÃO-HEURÍSTICA(estado)
        
    Se maximizando:
        max_eval = -INFINITO
        Para cada ação em GERAR-AÇÕES(estado):
            novo_estado = TRANSICAO(estado, ação)
            eval = MINIMAX(novo_estado, profundidade - 1, Falso)
            max_eval = MAX(max_eval, eval)
        Retorna max_eval
    Senão:
        min_eval = +INFINITO
        Para cada ação em GERAR-AÇÕES(estado):
            novo_estado = TRANSICAO(estado, ação)
            eval = MINIMAX(novo_estado, profundidade - 1, Verdadeiro)
            min_eval = MIN(min_eval, eval)
        Retorna min_eval
\end{verbatim}

\subsection{Alpha-Beta Pruning (Poda Alfa-Beta)}
A Poda Alfa-Beta é uma otimização obrigatória do Minimax que evita explorar ramos da árvore que não influenciarão 
a decisão final. Utiliza os parâmetros $\alpha$ (melhor escolha garantida para Max) e $\beta$ (melhor escolha garantida para Min).
Se em algum momento o algoritmo descobre que um estado pioraria a garantia já estabelecida ($\beta \leq \alpha$), 
o ramo inteiro é podado, ou seja, eliminado.

Isso reduz a complexidade para até $\mathcal{O}(b^{d/2})$ no melhor caso. 
Operando sob o mesmo limite fixo por jogada [cite: 15], o agente consegue atingir maiores profundidades.

\begin{verbatim}
Função ALPHA-BETA(estado, profundidade, alfa, beta, maximizando):
    Se profundidade == 0 ou TESTE-TERMINAL(estado):
        Retorna AVALIAÇÃO-HEURÍSTICA(estado)
        
    Se maximizando:
        max_eval = -INFINITO
        Para cada ação em GERAR-AÇÕES(estado):
            novo_estado = TRANSICAO(estado, ação)
            eval = ALPHA-BETA(novo_estado, profundidade - 1, alfa, beta, Falso)
            max_eval = MAX(max_eval, eval)
            alfa = MAX(alfa, eval)
            Se beta <= alfa:
                Interrompe o laço (PODA)
        Retorna max_eval
    ... (lógica espelhada para MIN)
\end{verbatim}

\subsection{Estratégia de Ordenação de Jogadas (Move Ordering)}
Para que a poda atinja seu desempenho máximo, implementou-se a ordenação de jogadas obrigatória. 
Lances que resultam em captura de peças adversárias recebem prioridade. Ao avaliar capturas primeiro, 
é muito provável que o algoritmo encontre rapidamente um valor elevado de $\alpha$, aumentando massivamente 
o número de podas prematuras (\textit{cut-offs}).

\begin{verbatim}
Função ORDENAR-JOGADAS(estado, acoes):
    Para cada ação em acoes:
        ação.score = 0
        Se TIPO-MOVIMENTO(ação) == CAPTURA:
            ação.score = 10
    Retorna acoes ordenadas de forma decrescente pelo score
\end{verbatim}

\end{document}